\begin{frame}
  \frametitle{Un código es un subespacio de qubits. Un código que corrige errores debe ser ortogonal y no ser deformado por el error}
    \begin{columns}[c]
      \begin{column}{0.5\textwidth}
        \centering
        % \resizebox ajusta el contenido al ancho de la línea actual (\linewidth)
        \resizebox{\linewidth}{!}{%
          \begin{tikzpicture}[>=Latex] % Better arrows for connection lines

            % Style for inner axes and vectors
            \tikzset{eje/.style={gray!70, thin, ->}}

            % Define colors for clarity
            \colorlet{contour}{black}
            \colorlet{inner}{gray!70}
            \colorlet{connection}{gray!50}

            % Two large container circles
            \node [draw, circle, minimum size=8.5cm, line width=0.8pt] (c1) at (0,0) {};
            \node [draw, circle, minimum size=8.5cm, line width=0.8pt] (c2) at (10,0) {};

            % Draw the central cube in the first circle
            \begin{scope}[shift={(c1.center)}]
              \coordinate (o_c) at (0,0);
              \coordinate (x_c) at (1, 0.2);
              \coordinate (y_c) at (-0.3, 1);
              \coordinate (z_c) at (-0.8, -0.5);

              % Outer contours with thicker lines
              \coordinate (xy_c) at ($ (x_c) + (y_c) $);
              \coordinate (xz_c) at ($ (x_c) + (z_c) $);
              \coordinate (yz_c) at ($ (y_c) + (z_c) $);
              \coordinate (xyz_c) at ($ (x_c) + (y_c) + (z_c) $);
              \draw[line width=0.8pt, contour] (x_c) -- (xy_c) -- (y_c);
              \draw[line width=0.8pt, contour] (y_c) -- (yz_c) -- (z_c);
              \draw[line width=0.8pt, contour] (z_c) -- (xz_c) -- (x_c);
              \draw[line width=0.8pt, contour] (xy_c) -- (xyz_c);
              \draw[line width=0.8pt, contour] (yz_c) -- (xyz_c);
              \draw[line width=0.8pt, contour] (xz_c) -- (xyz_c);

              % Inner axes with arrows, drawn after contours to be visible
              \draw[eje] (o_c) -- (x_c);
              \draw[eje] (o_c) -- (y_c);
              \draw[eje] (o_c) -- (z_c);
            \end{scope}

            % Draw parallelepipeds in the second circle
            % % A1: left side, vertically oriented -> Now a uniform cube
            \drawparal{7.5}{0.5}{(1,0.2)}{(-0.3,1)}{(-0.8,-0.5)}{$A_1$}{A1}

            % A2: upper-center, tall -> Now a uniform cube
            \drawparal{9.5}{2.5}{(1,0.2)}{(-0.3,1)}{(-0.8,-0.5)}{$A_2$}{A2}

            % A3: center, wide and regular -> Now a uniform cube
            \drawparal{10}{-0.5}{(1,0.2)}{(-0.3,1)}{(-0.8,-0.5)}{$A_3$}{A3}

            % A4: upper-right, flat -> Now a uniform cube
            \drawparal{12}{1.5}{(1,0.2)}{(-0.3,1)}{(-0.8,-0.5)}{$A_4$}{A4}

            % A5: lower-right, flat -> Now a uniform cube
            \drawparal{12}{-2}{(1,0.2)}{(-0.3,1)}{(-0.8,-0.5)}{$A_5$}{A5}

            % A6: lower-left, wide -> Now a uniform cube
            \drawparal{8}{-2.5}{(1,0.2)}{(-0.3,1)}{(-0.8,-0.5)}{$A_6$}{A6}
            % Connection lines from circle 1 center to circle 2 parallelepipeds
            \coordinate (orig) at (c1.center);

            % Draw connection lines with light grey and small arrow tips
            \draw [connection, thin, ->, shorten >=10pt, shorten <=10pt] (orig) -- (A1_orig);
            \draw [connection, thin, ->, shorten >=10pt, shorten <=10pt] (orig) -- (A2_orig);
            \draw [connection, thin, ->, shorten >=10pt, shorten <=10pt] (orig) -- (A3_orig);
            \draw [connection, thin, ->, shorten >=10pt, shorten <=10pt] (orig) -- (A4_orig);
            \draw [connection, thin, ->, shorten >=10pt, shorten <=10pt] (orig) -- (A5_orig);
            \draw [connection, thin, ->, shorten >=10pt, shorten <=10pt] (orig) -- (A6_orig);

          \end{tikzpicture}%
          }
        \end{column}

        % --- COLUMNA DERECHA: TEXTO EXPLICATIVO ---
        \begin{column}{0.5\textwidth}
          Un código cuántico es un subespacio del espacio de $n$ qubits (representado como $\mathcal{B}^{\otimes n}$ en esta presentación). Sin embargo, un código que corrige errores debe también cumplir con \citep{Kitaev2002}:
          \begin{enumerate}
            \item Ser ortogonal.
            \item Ser distinguible.
            \item No estar deformado.
          \end{enumerate}

          Con esto, un código $C$ corrige un error $\mathcal{E}$ si existe una operación $\mathcal{R}$ tal que:

          \begin{equation}
            \forall \rho \in C \implies (\mathcal{R} \circ \mathcal{E}) \rho \propto \rho
            \label{eq:error_correct}
          \end{equation}

        \end{column}
      \end{columns}
\end{frame}

\begin{frame}
  \frametitle{Las condiciones de Knill-Laflamme establecen que si un código ortogonal no se deforma con $\mathcal{E}$, entonces lo corrige.}

  Para un código cuántico ortogonal $C$ definido por el proyector $P$, este puede corregir las operaciones de un canal de ruido $\mathcal{E}(\rho) = \sum E_i \rho E_i^\dagger$ si \citep{Nielsen2010}:

    \begin{equation}
      P E_i^\dagger E_j P = \alpha_{ij} P
      \label{eq:knill-laflame}
    \end{equation}

  \vspace{0.5cm}
  Donde:
  \begin{itemize}
    \item \textbf{$\left\{E_i\right\}$}: son los operadores de Kraus del error $\mathcal{E}$.
    \item \textbf{$\alpha$}: es una matriz hermítica cuyas entradas $\alpha_{ij}$ \textbf{no dependen} del estado lógico codificado.
  \end{itemize}

  \textbf{Intuición:} Dado que $C$ es ortogonal, al no deformarse podemos seguir distinguiendo las palabras y corregir el error proyectando.

\end{frame}
